% Bsp. eines Hauptteils

\chapter{Introduction}
\label{sec:introduction}

\section{Motivation}
\label{sec:motivation}

Robotic grasping is a fundamental capability in automated manufacturing, logistics, and service robotics\cite{Todescato}. The ability to reliably grasp objects of
varying shapes, sizes, and surface properties remains a challenging problem, particularly in unstructured environments\cite{Dinakaran2022A}. Industrial applications such
as bin picking, assembly, and material handling require grippers that can operate robustly across a wide range of workpieces\cite{Kang2019Design}.

Among the commonly used end-effectors, suction-cup grippers and two-finger grippers represent two complementary approaches\cite{Courchesne2023A}. Suction grippers are
widely used for handling objects with smooth surfaces due to their simplicity and speed, whereas two-finger grippers provide flexibility in handling
objects with complex geometries. Designing optimal configurations for these grippers is non-trivial, as grasp performance depends on multiple
interacting factors such as geometry, contact distribution, force limits, and collision constraints\cite{Liu2024Topology}.

While the gripper configuration problems are well-motivated in themselves, the \textbf{primary objective of this thesis} is not simply to find the
best gripper configuration, but to \textbf{systematically apply and compare optimization algorithms} for solving these configuration problems. A large
number of optimization algorithms exist, spanning evolutionary strategies, decomposition-based methods, indicator-based selection, and Bayesian
optimization, yet their relative suitability for gripper design problems of this type is not well understood. Different algorithms may exhibit
substantially different convergence behaviour, solution quality, and robustness depending on the structure of the objective landscape and the
dimensionality of the search space. Identifying which algorithm class performs best for each gripper problem is therefore a practically relevant and
scientifically meaningful research question\cite{Doerterler2021Analyzing}.

To provide a rigorous and reproducible basis for this comparison, two gripper configuration problems are formulated as well-defined multi-objective
optimization tasks. In both cases, a diverse set of workpieces is incorporated directly into the problem formulation by treating each workpiece as a
separate objective. This design choice ensures that the optimized gripper configurations are not tailored to a single
geometry, but are instead robust across a range of workpiece shapes. It also increases the dimensionality of the objective space as the number of
workpieces grows, allowing the benchmarking study to probe algorithm behaviour under varying levels of problem complexity\cite{Guo2022Evolutionary}.

For the \textbf{multi-suction-cup gripper}, the objectives consist of the geometric stability of the suction base together with the grasp diversity
evaluated independently for each workpiece. These objectives are inherently conflicting: a larger suction base improves stability but increases the
risk of collisions with the workpiece, reducing the number of feasible grasps.

For the \textbf{two-finger gripper}, each workpiece contributes one objective, representing how well the optimized finger contour conforms to that
particular workpiece geometry. A finger contour that scores well across all workpieces simultaneously must generalize across diverse geometric
profiles, making the problem increasingly challenging as more workpieces are added\cite{Liu2024Topology}.

Together, these two problems provide a structured experimental basis for comparing optimization algorithms across varying objective space
dimensionalities and objective landscape characteristics, within the practically relevant domain of robotic gripper design.


% The \textbf{multi-suction-cup gripper} problem is cast as a bi-objective optimization task. The spatial arrangement of suction cups influences two
% competing performance aspects:

% \begin{enumerate} \item \textbf{Support Triangle Area:} The area of the geometric base formed by the suction cups determines grasp stability and
%     resistance to external disturbances. A larger base provides better force distribution and tilt resistance.

%     \item \textbf{Grasp Diversity:} The number, spatial location, and orientation of valid collision-free grasps achievable on a given workpiece. This
%     metric is inherently workpiece-dependent, as the gripper geometry must conform to the object's surface without collisions.
% \end{enumerate}

% These two objectives are inherently conflicting: a larger suction cup base improves stability but increases the likelihood of collisions with the
% workpiece, thereby reducing the number of feasible grasps. This conflict makes the problem a suitable benchmark for \textbf{multi-objective
% evolutionary algorithms} such as NSGA-II, NSGA-III, MOEA/D, SMS-EMOA, and CMOPSO, as well as \textbf{multi-objective Bayesian optimization} methods.

% The \textbf{two-finger gripper} problem is cast as a single-objective optimization task. The finger contour is parameterized by a spline defined
% through five points, three of which have variable vertical positions $h_{\mathrm{spline},i}$ ($i = 2, 3, 4$), along with the gripper nose length
% $l_{\mathrm{gn}}$ and depth $d_{\mathrm{gn}}$. The objective is to maximize the voxel-based proximity score $s_{\max}(\mathbf{q})$, computed via
% FFT-based convolution between the voxelized gripper and workpiece geometries. A configuration is valid only if the required finger opening
% corresponds to a physical closing motion. This five-dimensional, black-box, single-objective problem serves as a benchmark for
% \textbf{single-objective evolutionary algorithms} and \textbf{Bayesian optimization} methods.

\section{Problem Definition}
\label{sec:goals_tasks}
Designing grippers for a given set of workpieces is still largely a manual, experience-driven task, in which geometric parameters such as suction-cup
placement or finger contours are chosen iteratively and validated in simulation or on hardware\cite{Liu2022Simulation}. Because grasp quality can only
be assessed through a simulation or geometric evaluation routine, no analytical gradients or closed-form objective expressions are available, and each
evaluation is comparatively expensive\cite{Santoni2023Comparison}. Moreover, a single gripper usually has to handle several workpieces at once, so the
design criteria are inherently conflicting and cannot be reduced to a single scalar quality measure without losing information about the available
trade-offs. These properties of expensive, derivative-free, multi-objective evaluations place gripper design squarely in the domain of black-box
multi-objective optimization\cite{Mueller2017SOCEMO}. Which class of optimization algorithm is best suited for such problems, however, remains an open
question, as performance is known to depend strongly on the number of objectives, the shape of the objective landscape, and the available evaluation
budget\cite{Santoni2023Comparison, Raponi2023Optimizing}.

The primary goal of this thesis is to apply and comparatively evaluate optimization algorithms spanning evolutionary multi-objective methods and
Bayesian optimization for solving gripper configuration problems formulated as black-box multi-objective optimization
tasks\cite{Datta2011Optimizing,Belakaria2020Uncertainty-Aware}. The two gripper types, a multi-suction-cup gripper and a two-finger gripper, serve as
well-defined benchmark problems on which the algorithms are systematically compared with respect to solution quality, convergence behaviour, and
robustness.

To provide a rigorous and reproducible basis for this comparison, two gripper configuration problems are formulated as well-defined multi-objective
optimization tasks. In both cases, a diverse set of workpieces is incorporated directly into the problem formulation by treating each workpiece as a
separate objective. This design choice ensures that the optimized gripper configurations are not tailored to a single
geometry, but are instead robust across a range of workpiece shapes. It also increases the dimensionality of the objective space as the number of
workpieces grows, allowing the benchmarking study to probe algorithm behaviour under varying levels of problem complexity\cite{Guo2022Evolutionary}.

For the \textbf{multi-suction-cup gripper}, the objectives consist of the geometric stability of the suction base together with the grasp diversity
evaluated independently for each workpiece. These objectives are inherently conflicting: a larger suction base improves stability but increases the
risk of collisions with the workpiece, reducing the number of feasible grasps.

For the \textbf{two-finger gripper}, a fixed number of grasp points is defined for each workpiece, and the Objectives are computed from the conformance of the
finger contour to the corresponding workpiece section at each grasp point. A finger contour that scores well at all grasp points simultaneously must
therefore generalize across diverse geometric profiles, which makes the problem increasingly challenging as the number of grasp points grows.

Together, these two problems provide a structured experimental basis for comparing optimization algorithms across varying objective space
dimensionalities and objective landscape characteristics, within the practically relevant domain of robotic gripper design.

To achieve this goal, the following tasks are defined:

\begin{itemize}

    \item \textbf{Problem formulation:} Formulation of both gripper configuration problems as multi-objective optimization problems, including the
    definition of decision variables, objective functions, and constraints.
    
    \item \textbf{Optimization framework implementation:} Implementation of a generalized optimization framework capable of handling multi-objective
    problems of varying dimensionality and interfacing with both evaluation pipelines through a unified black-box interface.

    \item \textbf{Evaluation pipeline development:} Development of dedicated black-box evaluation pipelines for each gripper type, capable of
    assessing grasp quality given a gripper configuration and a target workpiece, and of serving as the objective function interface for all
    optimization algorithms.
    
    \item \textbf{Workpiece selection and experimental coverage:} Selection of a diverse set of workpieces spanning a range of geometric properties
    including varying surface curvatures, sizes, and shapes to ensure that algorithm comparisons are not biased towards a specific object geometry.
    Workpiece diversity is treated as a deliberate experimental design choice to expose differences in algorithm robustness and generalization across
    objective landscapes of varying complexity.

    \item \textbf{Algorithm application and comparison:} Application and systematic comparison of a selected set of multi-objective optimization
    algorithms drawn from the available literature \cite{pymoo, olson2025ax, botorch}, evaluated under controlled experimental conditions using
    well-defined performance indicators across both gripper types and workpieces\cite{Audet2018Performance}.

    \item \textbf{Result analysis:} Analysis and interpretation of algorithm performance across both gripper types and all workpieces, with particular
    focus on convergence behaviour, Pareto front quality, and the sensitivity of algorithm rankings to workpiece geometry and objective space
    dimensionality.
\end{itemize}


% ---------------------------------------------------------------
% End of chapter 
% ---------------------------------------------------------------

% ---------------------------------------------------------------
% End of chapter 
% ---------------------------------------------------------------
\chapter{Methodology}
\label{sec:methodology}

\section{Methodology for Multi-Suction Cup Gripper Optimization}

This section describes the optimization framework developed for a multi-suction
cup gripper consisting of three suction cups arranged in a triangular
configuration. The framework leverages the multi-objective optimization library
\textit{pymoo}~\cite{pymoo} and delegates objective function evaluation to an
external service via an API call, enabling a modular and decoupled architecture.

\subsection{Parameterization of the Gripper}

Each suction cup is defined by its planar position, while the height and
orientation remain fixed. Thus, the decision variables are:

\begin{equation}
    \mathbf{p} = [(x_1, y_1),\; (x_2, y_2),\; (x_3, y_3)]
\end{equation}
\nomenclature{$\mathbf{p}$}{Decision variable vector representing the gripper configuration}
\nomenclature{$(x_i, y_i)$}{Planar position of the $i$-th suction cup}

Each individual $\mathbf{p}$ forms a triangular base of the gripper.

\paragraph{Assumption}
A triangular configuration of suction cups is assumed to yield a stable grasp~\cite{madan2022realrobotchallenge2021}.
This is motivated by the fact that three spatially distributed contact points
form a statically stable support, enabling effective load distribution and
resistance to external disturbances.

\subsection{Decision Variables and Bounds}

Each decision variable is defined with a type (continuous or integer) and
associated bounds:

\begin{equation}
    x_i \in [\ell_i,\; u_i], \quad i = 1, \dots, n
\end{equation}
\nomenclature{$\ell_i$}{Lower bound of the $i$-th decision variable}
\nomenclature{$u_i$}{Upper bound of the $i$-th decision variable}
\nomenclature{$n$}{Number of decision variables}

The variable definitions are provided by the user as part of the optimization
request, allowing the framework to handle mixed-variable problems (i.e.,
combinations of continuous and integer variables).

\subsection{Optimization Framework}

The optimization is formulated as a multi-objective minimization problem and
solved using the \textit{pymoo} library. The user configures the optimization
through a request specifying:

\begin{itemize}
    \item Population size $N$
    \item Maximum number of generations $G_{\max}$
    \item Choice of optimization algorithm
    \item Number of objectives $n_{\text{obj}}$
    \item Decision variable definitions (type and bounds)
    \item Evaluation endpoint URL
\end{itemize}

\subsubsection{Supported Algorithms}

The framework supports the following multi-objective evolutionary algorithms:

\begin{table}[htbp]
\centering
\caption{Supported multi-objective optimization algorithms.}
\label{tab:algorithms}
\begin{tabular}{ll}
\toprule
\textbf{Identifier} & \textbf{Algorithm} \\
\midrule
NSGA2   & Non-dominated Sorting Genetic Algorithm II~\cite{nsga2} \\
NSGA3   & Non-dominated Sorting Genetic Algorithm III~\cite{nsga3} \\
MOEAD   & Multi-Objective Evolutionary Algorithm based on Decomposition~\cite{moead} \\
CMOPSO  & Competitive Multi-Objective Particle Swarm Optimization \\
SMSEMOA & S-Metric Selection Evolutionary Multi-Objective Algorithm~\cite{smsemoa} \\
\bottomrule
\end{tabular}
\end{table}

For algorithms requiring reference directions (NSGA-III, MOEA/D), uniform
reference directions are generated using Das and Dennis's
method~\cite{das_dennis}:

\begin{equation}
    \Lambda = \texttt{get\_reference\_directions}(\text{``uniform''},\; n_{\text{obj}},\; n_{\text{partitions}})
\end{equation}
\nomenclature{$\Lambda$}{Set of reference directions for decomposition-based algorithms}

For mixed-variable problems, specialized sampling and mating operators are
employed to handle heterogeneous variable types.

\subsection{Initialization}

An initial population of $N$ candidate solutions is generated via mixed-variable
sampling:

\begin{equation}
    \mathcal{P}^{(0)} = \{\mathbf{p}_1, \mathbf{p}_2, \dots, \mathbf{p}_N\}
\end{equation}
\nomenclature{$\mathcal{P}^{(0)}$}{Initial population of candidate solutions}
\nomenclature{$N$}{Population size}

\subsection{Objective Function Evaluation via API}

The evaluation of objective function values is decoupled from the optimization
loop and delegated to an external evaluation service. For each candidate
solution $\mathbf{p}$, the optimizer issues an API call to the evaluation
endpoint, which executes the grasp evaluation subprocess (described in
Section~\ref{sec:grasp_eval}) and returns the objective values.

\begin{equation}
    \mathbf{f}(\mathbf{p}) = \texttt{API\_call}(\mathbf{p}) \rightarrow [f_1(\mathbf{p}),\; f_2(\mathbf{p})]
\end{equation}

This architecture enables:
\begin{itemize}
    \item Separation of concerns between optimization logic and domain-specific evaluation
    \item Scalability through independent deployment of the evaluation service
    \item Reusability of the optimization framework for different problem domains
\end{itemize}

\subsubsection{Grasp Evaluation Subprocess}
\label{sec:grasp_eval}

The evaluation service receives a candidate gripper configuration $\mathbf{p}$
and performs the following steps:

\begin{enumerate}

\item \textbf{Grasp Generation:}
For each point in the workpiece point cloud, a candidate suction grasp is
generated by aligning the tool center point (TCP) of the gripper with the
surface point.

\item \textbf{Collision Filtering:}
The gripper is modeled as a rigid 3D body with fixed height, forming a
triangular prism whose base is defined by $\mathbf{p}$. Each candidate grasp is
checked for collisions between the gripper body and the workpiece. Colliding
grasps are discarded.

\item \textbf{Clustering of Valid Grasps:}
Let $\mathcal{G}_{\text{valid}}$ denote the set of collision-free grasps. If
$|\mathcal{G}_{\text{valid}}| \geq K_{\min}$, clustering is performed based on
spatial proximity and orientation similarity to reduce redundancy.
\nomenclature{$\mathcal{G}_{\text{valid}}$}{Set of valid (collision-free) grasps}
\nomenclature{$K_{\min}$}{Minimum number of valid grasps required for clustering}

\item \textbf{Objective Computation:}
The service computes and returns the objective values:

\begin{equation}
    \mathbf{f}(\mathbf{p}) = [-D(\mathbf{p}),\; -A(\mathbf{p})]
\end{equation}

where:
\begin{itemize}
    \item $A(\mathbf{p})$ is the triangle area (maximized for stability):
    \begin{equation}
        A(\mathbf{p}) = \frac{1}{2} \left| (x_2 - x_1)(y_3 - y_1) - (x_3 - x_1)(y_2 - y_1) \right|
    \end{equation}
    \item $D(\mathbf{p})$ is the diversity metric (maximized for robustness),
    computed from the spatial distribution and orientation spread of clustered grasps.
\end{itemize}

\nomenclature{$A(\mathbf{p})$}{Area of the triangle formed by the three suction cups}
\nomenclature{$D(\mathbf{p})$}{Diversity metric of the grasp distribution}
\nomenclature{$\mathbf{f}(\mathbf{p})$}{Multi-objective function vector}

Since \textit{pymoo} minimizes by convention, maximization objectives are
negated.

\end{enumerate}

\subsection{Optimization Loop}

The evolutionary optimization proceeds as follows:

\begin{equation}
    \mathcal{P}^{(t+1)} = \mathcal{A}\left(\mathcal{P}^{(t)},\; \mathbf{f}\right), \quad t = 0, 1, \dots
\end{equation}
\nomenclature{$\mathcal{A}$}{Selected multi-objective evolutionary algorithm}
\nomenclature{$t$}{Generation index}

where $\mathcal{A}$ denotes the selected algorithm performing selection,
crossover, mutation, and survival operations.

\subsection{Termination Criteria}

The optimization terminates when any of the following conditions is satisfied:

\begin{table}[htbp]
\centering
\caption{Termination criteria.}
\label{tab:termination}
\begin{tabular}{lll}
\toprule
\textbf{Criterion} & \textbf{Symbol} & \textbf{Description} \\
\midrule
Design space tolerance  & $x_{\text{tol}}$ & Change in decision variables below threshold \\
Objective space tolerance & $f_{\text{tol}}$ & Relative improvement in objectives below threshold \\
Maximum generations     & $G_{\max}$       & Maximum number of generations reached \\
Maximum evaluations     & $E_{\max}$       & Maximum number of function evaluations reached \\
\bottomrule
\end{tabular}
\end{table}
\nomenclature{$x_{\text{tol}}$}{Design space convergence tolerance}
\nomenclature{$f_{\text{tol}}$}{Objective space convergence tolerance}
\nomenclature{$G_{\max}$}{Maximum number of generations}
\nomenclature{$E_{\max}$}{Maximum number of function evaluations}

Tolerances are evaluated over a sliding window of $w$ generations.

\subsection{Performance Monitoring}

The hypervolume indicator~\cite{hypervolume} is tracked at each generation via a
callback mechanism:

\begin{equation}
    \text{HV}^{(t)} = \text{HV}(\mathcal{F}^{(t)},\; \mathbf{r})
\end{equation}
\nomenclature{$\text{HV}^{(t)}$}{Hypervolume indicator at generation $t$}
\nomenclature{$\mathbf{r}$}{Reference point for hypervolume computation}
\nomenclature{$\mathcal{F}^{(t)}$}{Objective values of the population at generation $t$}

where $\mathbf{r} \in \mathbb{R}^{n_{\text{obj}}}$ is the reference point and
$\mathcal{F}^{(t)}$ denotes the set of objective vectors at generation $t$. A
monotonically increasing hypervolume indicates convergence toward the Pareto
front.

\subsection{Optional Hyperparameter Optimization}

The framework optionally supports automated tuning of algorithm hyperparameters
(e.g., crossover probability, mutation rates) using a meta-optimization
approach. A \texttt{HyperparameterProblem} is constructed around the selected
algorithm and evaluated over multiple random seeds. The mean hypervolume across
runs serves as the meta-objective:

\begin{equation}
    \min_{\boldsymbol{\theta}} \left\{ -\frac{1}{|\mathcal{S}|} \sum_{s \in \mathcal{S}} \text{HV}(\mathcal{F}_s,\; \mathbf{r}) \right\}
\end{equation}
\nomenclature{$\boldsymbol{\theta}$}{Algorithm hyperparameters}
\nomenclature{$\mathcal{S}$}{Set of random seeds for meta-optimization}

The best hyperparameters are then applied to the main optimization run.

\subsection{Output}

The optimization produces a set of Pareto-optimal solutions:

\begin{equation}
    \mathcal{P}^* = \left\{ \mathbf{p} \;\middle|\; \nexists\; \mathbf{p}' \in \mathcal{P} : \mathbf{f}(\mathbf{p}') \preceq \mathbf{f}(\mathbf{p}) \right\}
\end{equation}
\nomenclature{$\mathcal{P}^*$}{Set of Pareto-optimal solutions}

The final output includes:
\begin{itemize}
    \item The Pareto-optimal set $\mathcal{P}^*$ with corresponding objective values
    \item The hypervolume progression $\{\text{HV}^{(0)}, \text{HV}^{(1)}, \dots, \text{HV}^{(T)}\}$
    \item Total optimization wall-clock time
\end{itemize}

These solutions represent optimal trade-offs between grasp diversity and
triangle area (stability).

\chapter{Stand der Technik}
\label{sec:real-unter}
Beschreibung des momentanen Standes der Technik und Wissenschaft. Hier muss auch
klar definiert werden, in welchem Zusammenhang das Problem steht, d.h.
der Systemkontext und der bekannte L�sungsraum ist klar zu beschreiben 
(vorhandene.L�sungen

\newpage
\section{Problembeschreibung}
\label{sec:real-unterWeiter}

\newpage
\section{Existierende L�sungen}
\label{sec:real-unterWeiterNoch}


\section{Literaturverweise}
\label{sec:real-literatur}

Verweise im Text: \cite{doc:stz} und \cite{doc:gun}.

\chapter{L�sungsansatz}
\label{sec:loesung}
Hier k�nnen je nach Art der Arbeit auch verschiedene L�sungsans�tze
beschrieben und diskutiert werden.

\chapter{Durchgef�hrte Arbeiten}
\label{sec:erledigt}
Es muss klar erkennbar sein, auf was aufgesetzt wurde, d.h. schon vorhanden
war, und was im Rahmen dieser Arbeit selbst entwickelt wurde.

\chapter{Ergebnisse}
\label{sec:ergebnis}
Hier kann z.B. eine Bedienanleitung stehen, oder bei Analysen die Analyseergebnisse.
\enquote{Neuigkeiten} Messergebnisse

\chapter{Zusammenfassung}
\label{sec:zus}
Kurze Zusammenfassung und evtl. Beschreibung der Punkte, die zu einer
kompletten, marktreifen L�sung noch fehlen.


